\section{Modules and Long-Range Interactions}
\label{sec:intro-modules}

\subsection{Random Notes on 2D Architectures}

The main point should be shifted, in the intro, and maybe elucidated here,
that low-depth for factoring can be achieved by moving into 2D, but
moreover that the width and size blowups can also be reduced with a
hybrid 2D module.

The model is hybrid, because instead of a single, contiguous 2D lattice,
say with $W$ qubits and $S$ interactions between qubits (counting the
general case of all single- and two-qubit gates as interactions),
the computation is split among many communicating 2D lattices
(say $\overline{W}$ of them), each containing $W/\overline{W}$ qubits
each. In this new model, we argue that there are now a total number of
qubit interactions $S'$ across all modules, where $S' \le S$, and at
the same time we have introduced $\overline{S}$
new, \emph{long-range} interactions between modules. Based on experimental
data of create shared entanglement between separate ion traps (citation needed here),
these long-range interactions are more expensive than short-range interactions
within a module, and so we count them separately.

Currently, we impose no constraints on the connectivity of the modules, which exist
on a higher-level graph of modules. This is supported by current proposals
(cite MUSIQC here). However, it is interesting to see what effect such constraints
would have on a factoring architecture. We will return to this later.

With no such constraints (arbitrary connectivity of modules), we can now see the
effect of setting the size of each module. If we set modules
to be a single qubit in size ($\overline{W} = W$), we would get the architectural
model \textsf{AC}. If we were to set module size to include the entire lattice
($\overline{W} = 1$), then we would get the model described in the previous
section, \textsf{2D CCNTC}. This new model then, \textsf{2D CCNTCM}, is a hybrid
between \textsf{AC} and \textsf{2D CCNTC} with module size as a parameter.
\textsf{AC} and \textsf{2D CCNTC} represent the two extremes of long-range
interactions, where \textsf{AC} permits all long-range interactions and
\textsf{2D CCNTC} permits none.

There are now two more things to study and discuss. One is the optimal setting
of the module size $W/\overline{W}$, which in this paper we have set to
be $O(n)$. The second is the connectivity constraints of the modules.

For the first parameter, we have not rigorously shown that the optimal setting
of $W / \overline{W}$ is 

Now then, what is the effect of imposing constraints on module connectivity?

\subsection{\textsf{2D CCNTCM}: Two-Dimensional Nearest-Neighbor Two-Qubit Concurrent Gates with Classical Controller and Modules}
\label{subsec:2dccntcm}

A single, contiguous
2D lattice that contains an entire quantum architecture may be prohibitively large to manufacture. In practice,
scalable experiments will probably use many
smaller quantum computers that communicate by means of shared
entanglement \cite{Monroe2012}.
We call these individual machines \emph{modules}, each of
which is a self-contained \textsf{2D CCNTC} lattice. This should not be
confused with the word ``modular'' as in ``modular arithmetic'' or as
referring to the modulus $m$ that we are trying to factor.

We treat these modules
and teleportations between them as nodes and edges, respectively,
in a higher-level planar graph. The teleportations each transmit one qubit
from one module to another, from any location within the source module
to any location within the destination module, making use of the
omnipresent classical controller. The modules can be arbitrarily far
apart physically, but they have bounded-degree connectivity with other
modules, and their edges are planar (they cannot intersect).

A single module can be part of multiple teleportation operations in a single timestep, as long as they involve disjoint qubits within the module.
We justify this assumption in that it is
possible to establish entanglement between multiple
quantum computers
in parallel. We call this new model \textsf{2D CCNTCM},
and we argue that is captures the essential aspects of 2D architectures
without being overly sensitive to the exact geometry of the lattices involved.
A graphic depiction of three modules in \textsf{2D CCNTCM} is shown in
Figure \ref{fig:modules}. Each module contains within it a
\textsf{2D CCNTC} lattice. We can equivalently consider the omnipresent,
single
classical controller as a collection of multiple classical controllers, one
for each module or teleportation operation, which can inter-communicate
classically and share a clock.

\textsf{2D CCNTCM} is very similar to the model of a quantum multicomputer
of Van Meter
\cite{VanMeter2006} with the difference that we allow parallel teleportations.
It is also very similar to the distributed quantum computer \textsf{DQC} of
Beals et al. \cite{Beals2012} except that we require each module to have
$\Omega(n)$ qubits in order to enforce some locality.

\begin{figure}[btp!]
\begin{center}
\includegraphics[width=4in]{factor-polylog/figures/modules.pdf}
\end{center}
\caption{Three modules in the \textsf{2D CCNTCM} model}
\label{fig:modules}
\end{figure}

\begin{definition}
A \textsf{2D CCNTCM} architecture consists of

\begin{itemize}
\item a quantum computer $\overline{QC}$ which is represented by a planar graph $(\overline{V},\overline{E})$. A
node $\overline{v} \in \overline{V}$ represents a module, or a graph $(V,E)$
from a \textsf{2D CCNTC} architecture defined previously. It can have
unbounded degree.
An
undirected edge $(\overline{u},\overline{v}) \in \overline{E}$ represents an
allowed teleportation from any qubit in module $\overline{u}$ to
another qubit in module $\overline{v}$.
\item All modules are restricted to be linear in the number of their qubits:
$|V| = \Theta(n)$ for all $(V,E) \in \overline{V}$.
\item a universal gate set $\mathcal{G} = \{X, Z, H, T, T^{\dagger}, CNOT,
MeasureZ\}$
for the qubits \emph{within the same} modules which is the same as for \textsf{2D CCNTC},
and an additional operation $Teleport$ which only operates on qubits
\emph{in
different} modules.
\item a deterministic machine (classical controller) $\overline{CC}$ that applies a sequence
of concurrent gates in each of $D+\overline{D}$ timesteps.
This can be a separate classical controller
for every pair of modules.
\item In timestep $i$, $\overline{CC}$ applies
gates $G_i = \{g_{i,j} : g_{i,j} \in \mathcal{G} \lor g_{i,j} = Teleport \}$.
That is, there are two kinds of timesteps with respect to the kinds of gates
which operate within them.
\begin{enumerate}
\item In the first kind, gates are exclusively from $\mathcal{G}$, and
they operate within modules as described
for \textsf{2D CCNTC} above. We say there are $D$ such timesteps.
\item In the second kind, gates are exclusively $Teleport$ gates between two qubits $v^{(1)}_{i,j} \in \overline{v}_1$ and
$v^{(2)}_{i,j} \in \overline{v}_2$ for
(possibly non-distinct) modules $\overline{v}_1, \overline{v}_2 \in \overline{V}$.
Again, all such qubits much be distinct within a timestep.
We say there are $\overline{D}$ such timesteps.
\end{enumerate}

Again, we define the support of $G_i$
as $V_i$, the set of all qubits acted upon by any $g_{i.j}$, which
includes all the modules.
\begin{equation}
V_i = \bigcup_{j: g_{i,j} \in G_i} v_{i,j} \cup v^{(1)}_{i,j} \cup v^{(2)}_{i,j} 
\end{equation}

\end{itemize}
\end{definition}

We measure the efficiency of a circuit in this new module using not just
the three conventional circuit resources, but with three novel resources
based on modules.

\begin{description}

%, depicted in Figure \ref{fig:resources}:
\item[module depth ($\overline{D}$):] the depth of consecutive teleportations between modules.
\item[module size ($\overline{S}$):] the number of total qubits teleported between any two modules over all timesteps.
\item[module width ($\overline{W}$):] the number of modules whose qubits are
acted upon during any timestep.

\end{description}

%We can make an observation analogous to Equation \ref{eqn:depth-width} but
%for modules in Equation \ref{eqn:module-depth-width}.

%\begin{equation}
%\overline{D} \le \overline{S} \le \overline{D}\cdot \overline{W}
%\label{eqn:module-depth-width}
%\end{equation}

We note the following relationship between circuit width and
module width.

\begin{equation}
W = O(n\overline{W})
\label{eqn:module-width}
\end{equation}

This restriction imposes some locality on our model by constraining it to
nearest-neighbor gates within a linear-sized group of qubits, but allowing
it long-range teleportation to circumvent onerous geometric constraints.
Using the constant-depth communication in Section \ref{sec:cdc}, and for
the specific case of factoring, we
can simulate arbitrary connectivity between modules with only a polynomial
increase in the module size and a constant increase in module depth.
